\chapter{Password Hashing with bcrypt}

A breach of a database storing plain-text passwords compromises every account instantly -- and, since people reuse passwords, their accounts elsewhere too. This chapter's rule: passwords are never stored as-is.

\begin{warnbox}[WARNING]
Never store passwords in plain text or with reversible encryption -- a leaked encryption key exposes every password. Hashing is one-way, which is why bcrypt (not AES or a custom cipher) is the right tool.
\end{warnbox}

\section{Hashing and Salt}

A hash function maps any input to a fixed-length, one-way output -- the same input always yields the same hash. A plain hash (like SHA-256) gives identical hashes for identical passwords, letting attackers crack many accounts at once with precomputed ``rainbow tables.''

A \textbf{salt} -- random data mixed in before hashing -- makes identical passwords hash differently. bcrypt generates and stores the salt automatically inside the hash string, so you never manage it manually.

\begin{diagrambox}[bcrypt Hashing Flow]
\begin{verbatim}
 "MySecret123" -> bcrypt generates a random salt
        |
        v
 salt + password -> bcrypt algorithm (2^rounds iterations)
        |
        v
 stored hash: $2b$10$N9qo8uLOickgx2ZMRZoMy...
 (algorithm + rounds + salt + hash, all in one string)
\end{verbatim}
\end{diagrambox}

\section{Hashing on Registration}

Install bcrypt and hash the password before saving the user document.

\begin{lstlisting}[language=JavaScript]
// controllers/authController.js
import bcrypt from "bcrypt";
import User from "../models/User.js";

export const registerUser = async (req, res, next) => {
  try {
    const { name, email, password } = req.body;

    const existing = await User.findOne({ email });
    if (existing) {
      return res.status(409).json({ success: false, message: "Email already in use" });
    }

    const salt = 10; // salt rounds
    const hashedPassword = await bcrypt.hash(password, salt);

    const user = await User.create({
      name,
      email,
      password: hashedPassword,
      role: "user",
    });

    res.status(201).json({
      success: true,
      data: { id: user._id, name: user.name, email: user.email, role: user.role },
    });
  } catch (err) {
    next(err);
  }
};
\end{lstlisting}

\subsection{The Salt Rounds Tradeoff}

The second argument to \texttt{bcrypt.hash} is the \textbf{cost factor} (salt rounds) -- it controls the internal loop count, and work scales as $2^{\text{rounds}}$.

\begin{center}
\begin{tabularx}{\textwidth}{lXX}
\toprule
\textbf{Rounds} & \textbf{Effect} & \textbf{Tradeoff} \\
\midrule
8--10 & Fast (ms to ~100ms) & Weaker, fine for most apps \\
12 & Slower (~250ms) & Better resistance, adds latency \\
14+ & Very slow (seconds) & Strong, hurts UX/throughput \\
\bottomrule
\end{tabularx}
\end{center}

\begin{notebox}[NOTE]
10 rounds is the standard default -- roughly 50--100ms per hash, imperceptible to users but expensive enough to deter brute forcing.
\end{notebox}

\section{Verifying on Login}

You never ``unhash'' a password -- bcrypt re-hashes the submitted password using the salt embedded in the stored hash and compares the results.

\begin{lstlisting}[language=JavaScript]
// controllers/authController.js
export const loginUser = async (req, res, next) => {
  try {
    const { email, password } = req.body;

    const user = await User.findOne({ email });
    if (!user) {
      return res.status(401).json({ success: false, message: "Invalid credentials" });
    }

    const match = await bcrypt.compare(password, user.password);
    if (!match) {
      return res.status(401).json({ success: false, message: "Invalid credentials" });
    }

    // generateToken + res.cookie happen here (see Chapters 18-19)
    res.status(200).json({ success: true, data: { id: user._id, email: user.email } });
  } catch (err) {
    next(err);
  }
};
\end{lstlisting}

\begin{warnbox}[WARNING]
Return the same generic ``Invalid credentials'' message whether the email doesn't exist or the password is wrong -- distinguishing the two leaks which emails are registered.
\end{warnbox}

\section{Alternative: Hashing in a Mongoose Pre-Save Hook}

Centralize hashing in the User schema with a \texttt{pre("save")} hook instead of repeating it in every controller -- hashing then happens no matter where the document is saved from.

\begin{lstlisting}[language=JavaScript]
// models/User.js
import mongoose from "mongoose";
import bcrypt from "bcrypt";

const userSchema = new mongoose.Schema({
  name: { type: String, required: true },
  email: { type: String, required: true, unique: true },
  password: { type: String, required: true },
  avatar: { type: String, default: "" },
  role: { type: String, enum: ["user", "manager", "admin"], default: "user" },
});

userSchema.pre("save", async function (next) {
  if (!this.isModified("password")) return next(); // skip if unchanged

  const salt = await bcrypt.genSalt(10);
  this.password = await bcrypt.hash(this.password, salt);
  next();
});

userSchema.methods.comparePassword = function (candidate) {
  return bcrypt.compare(candidate, this.password);
};

export default mongoose.model("User", userSchema);
\end{lstlisting}

\subsection{Why \texttt{isModified("password")} Matters}

Without this check, saving unrelated fields (like \texttt{avatar}) would re-hash the already-hashed password, locking the user out. \texttt{isModified("password")} ensures re-hashing only happens when the raw password field actually changed.

\begin{tipbox}[TIP]
With this approach the controller simplifies to \texttt{User.create(\{ name, email, password, role \})} -- login still uses \texttt{user.comparePassword(password)}.
\end{tipbox}

\whatwhyhow
{bcrypt hashes passwords with a per-password random salt before storage, and compares candidate passwords against the stored hash on login.}
{Plain-text or reversibly-encrypted passwords turn any database leak into a full account compromise; one-way hashing with salt makes stored values useless to an attacker even if stolen.}
{Hash with \texttt{bcrypt.hash(password, 10)} on registration (or via a \texttt{pre("save")} hook guarded by \texttt{isModified("password")}), and verify with \texttt{bcrypt.compare(password, user.password)} on login.}
